\section{Подписанные данные}\label{FMT.SignedData}

Формат подписанных данных задается типом \code{SignedData}, который определен в
СТБ~34.101.23.

Контейнер~\code{SignedData} заполняется по правилам СТБ~34.101.23 со следующими
уточнениями.

\begin{enumerate}
\item
Версия синтаксиса (компонент \code{version}) должна равняться~$3$.

\item
Список идентификаторов алгоритмов хэширования (компонент
\code{digestAlgorithms}) должен содержать единственный элемент, и этот элемент
должен быть выбран из перечня, заданного в~\ref{CRYPTO.Hash}.

\item
Тип подписываемых данных (компонент~\code{eContentType}, вложенный 
в~\code{encapContentInfo}) должен принимать одно из следующих значений:

\begin{enumerate}
\item
\code{id-ct-TSTInfo}, если подписывается ответ СШВ (см.~\ref{FMT.TSP.Resp});
\item
\code{id-ct-DVCSResponseData}, если подписывается ответ СЗД
(см.~\ref{FMT.DVCS.Resp});
\item
\code{bpki-ct-enroll1-req}, если в сценарии~\code{Enroll1}
процесса~\code{Enroll} подписывается запрос на выпуск сертификата
(см.~\ref{PROCESSES.Enroll.Signed});
\item
\code{bpki-ct-enroll2-req}, если в сценарии~\code{Enroll2}
процесса~\code{Enroll} подписывается запрос на выпуск сертификата
(см.~\ref{PROCESSES.Enroll.Signed});
\item
\code{bpki-ct-reenroll-req}, если в процессе~\code{Reenroll} подписывается
запрос на выпуск сертификата (см.~\ref{PROCESSES.Reenroll});
\item
\code{bpki-ct-spawn-req}, если в процессе~\code{Spawn} подписывается запрос на
выпуск сертификата (см.~\ref{PROCESSES.Spawn});
\item
\code{bpki-ct-setpwd-req}, если в процессе~\code{Setpwd} подписывается новый
пароль (см.~\ref{PROCESSES.Setpwd});
\item
\code{bpki-ct-revoke-req}, если в процессе~\code{Revoke} подписывается запрос на
отзыв сертификата (см.~\ref{PROCESSES.Revoke});
\item
\code{bpki-ct-resp}, если подписывается ответ УЦ (см.~\ref{FMT.BPKIResp}).
\end{enumerate}

Первый идентификатор определен в СТБ~34.101.82, второй~--- в СТБ~
34.101.81, остальные~--- в приложении~\ref{ASN1}.

\item
Опциональный компонент~\code{certificates} должен присутствовать и 
должен содержать единственный сертификат~--- сертификат подписанта.
%
Опциональный компонент~\code{crls} должен быть опущен.

\item
Список данных о подписантах (компонент~\code{signerInfos}) должен содержать
единственный контейнер типа~\code{SignerInfo}, и этот контейнер должен быть
заполнен следующим образом:
\begin{enumerate}
\item
версия (компонент \code{version}) должна равняться~$1$;
\item
идентификатор подписанта (\code{sid}) должен быть задан через
тип~\code{IssuerAndSerialNumber}. Компоненты~\code{issuer} и~\code{serialNumber}
этого типа должны повторять одноименные компоненты сертификата подписанта;
\item
идентификатор алгоритма хэширования (\code{digestAlgorithm}) должен совпадать с
идентификатором, указанным в компоненте~\code{digestAlgorithms} основного
контейнера~\code{SignedData};
\item
идентификатор алгоритмов ЭЦП (\code{signatureAlgorithm}) должен быть выбран из
перечня, заданного в~\ref{CRYPTO.Sign}. Алгоритмы ЭЦП должны соответствовать
алгоритму хэширования из~\code{digestAlgorithm} и открытому ключу подписанта;
\item
в список подписанных атрибутов (\code{signedAttrs}) должны быть включены
атрибуты <<Тип содержимого>> (\code{ContentType}), <<Хэш-значение>>
(\code{MessageDigest}) и может быть включен атрибут <<Время подписания>>
(\code{SigningTime}). Все атрибуты определены в СТБ~34.101.23.
%
Списки подписанных атрибутов штампа времени и аттестата заверения должны
дополнительно включать атрибут \code{SigningCertificateV2}, определенный в
СТБ~34.101.80. В этом атрибуте должна быть указана ссылка на сертификат
подписанта, ссылка должна быть сформирована с помощью алгоритма \code{belt-hash}
и должна быть единственной;
\item
список неподписанных атрибутов (\code{unsignedAttrs}) должен быть пуст.
\end{enumerate}
\end{enumerate}
